\documentclass[12pt,oneside]{book}

% ============================================================
% METADATOS
% ============================================================
\newcommand{\universidad}{Universidad de Buenos Aires (UBA)}
\newcommand{\facultad}{Facultad de Derecho}
\newcommand{\carrera}{Abogacía}
\newcommand{\materia}{Derecho Procesal Civil y Comercial}
\newcommand{\titulo}{Recurso de Apelación}
\newcommand{\autor}{Isaac Edgar Camacho Ocampo}
\newcommand{\anio}{2024}

% ============================================================
% PAQUETES
% ============================================================
\usepackage[utf8]{inputenc}
\usepackage[T1]{fontenc}
\usepackage[spanish,es-nodecimaldot]{babel}

\usepackage[
    top=2.5cm,
    bottom=2.5cm,
    left=3cm,
    right=2.5cm,
    headheight=15pt
]{geometry}

\usepackage{amsmath}
\usepackage{amssymb}
\usepackage{mathtools}

\usepackage{graphicx}
\usepackage{float}
\usepackage{caption}
\usepackage{subcaption}

\usepackage{booktabs}
\usepackage{array}
\usepackage{longtable}
\usepackage{tabularx}

\usepackage{enumitem}

\usepackage{xcolor}
\usepackage[most]{tcolorbox}

\usepackage{fancyhdr}

\usepackage{ragged2e}

% ============================================================
% RUTAS DE IMÁGENES
% ============================================================
\graphicspath{
    {./img/}
    {./graficos/}
    {./figuras/}
}

% ============================================================
% CONFIGURACIÓN
% ============================================================
\setcounter{tocdepth}{2}

\setlength{\parindent}{0pt}
\setlength{\parskip}{0.5em}

\setlist[itemize]{
    topsep=0.3em,
    itemsep=0.2em
}

\setlist[enumerate]{
    topsep=0.3em,
    itemsep=0.2em
}

\clubpenalty=10000
\widowpenalty=10000

\pagestyle{fancy}
\fancyhf{}
\fancyhead[L]{\nouppercase{\leftmark}}
\fancyhead[R]{\materia}
\fancyfoot[C]{\thepage}
\renewcommand{\headrulewidth}{0.4pt}
\renewcommand{\footrulewidth}{0pt}

\fancypagestyle{plain}{
    \fancyhf{}
    \fancyfoot[C]{\thepage}
    \renewcommand{\headrulewidth}{0pt}
}

% ============================================================
% COMANDOS MATEMÁTICOS
% ============================================================
\newcommand{\abs}[1]{\left|#1\right|}
\newcommand{\norm}[1]{\left\lVert#1\right\rVert}
\newcommand{\vect}[1]{\mathbf{#1}}
\newcommand{\deriv}[2]{\frac{d#1}{d#2}}
\newcommand{\pderiv}[2]{\frac{\partial#1}{\partial#2}}

% ============================================================
% ENTORNOS / CAJAS ACADÉMICAS
% ============================================================
\newtcolorbox{definition}[1]{
    enhanced,
    breakable,
    title={Definición: #1},
    fonttitle=\bfseries,
    colback=blue!3,
    colframe=blue!50!black,
    boxrule=0.8pt,
    arc=2mm
}

\newtcolorbox{important}{
    enhanced,
    breakable,
    title={Importante},
    fonttitle=\bfseries,
    colback=yellow!8,
    colframe=orange!70!black,
    boxrule=0.8pt,
    arc=2mm
}

\newtcolorbox{example}[1]{
    enhanced,
    breakable,
    title={Ejemplo: #1},
    fonttitle=\bfseries,
    colback=green!4,
    colframe=green!40!black,
    boxrule=0.8pt,
    arc=2mm
}

\newtcolorbox{formula}[1]{
    enhanced,
    breakable,
    title={Fórmula / Regla: #1},
    fonttitle=\bfseries,
    colback=purple!4,
    colframe=purple!50!black,
    boxrule=0.8pt,
    arc=2mm
}

\newtcolorbox{exercise}[1]{
    enhanced,
    breakable,
    title={Ejercicio: #1},
    fonttitle=\bfseries,
    colback=gray!5,
    colframe=gray!60!black,
    boxrule=0.8pt,
    arc=2mm
}

\newtcolorbox{note}{
    enhanced,
    breakable,
    title={Nota},
    fonttitle=\bfseries,
    colback=white,
    colframe=gray!50,
    boxrule=0.6pt,
    arc=2mm
}

% ============================================================
% HIPERVÍNCULOS Y METADATOS PDF
% ============================================================
\usepackage[
    colorlinks=true,
    linkcolor=blue!50!black,
    citecolor=green!40!black,
    urlcolor=blue!60!black,
    pdftitle={\titulo},
    pdfauthor={\autor},
    pdfsubject={Apuntes universitarios},
    bookmarks=true
]{hyperref}

% ============================================================
% DOCUMENTO
% ============================================================
\begin{document}

% ---------- PORTADA ----------
\begin{titlepage}
\thispagestyle{empty}

\begin{center}

\vspace*{1cm}

{\large \textsc{\universidad}}

\vspace{0.2cm}

{\large \textsc{\facultad}}

\vspace{0.2cm}

{\large \carrera}

\vspace{3cm}

{\Huge \textbf{\materia}}

\vspace{1cm}

{\LARGE \titulo}

\vspace{0.8cm}

\textit{Comprender los fundamentos de cada instituto procesal es el primer paso para transformar el estudio en criterio jurídico.}

\vspace{1.2cm}

\rule{0.70\textwidth}{0.5pt}

\vspace{0.5cm}

{\large Apuntes de estudio}

\vspace{0.5cm}

\rule{0.70\textwidth}{0.5pt}

\vfill

{\large \autor}

\vspace{0.3cm}

{\large \anio}

\end{center}

\vfill

\begin{flushleft}
\textit{Este es un modesto aporte a los estudiantes de ``Derecho Procesal Civil y Comercial'' no pretende sustituir el material oficial de la materia.}
\end{flushleft}

\end{titlepage}

% ---------- ÍNDICE ----------
\tableofcontents

% ============================================================
% CONTENIDO
% ============================================================

\chapter{Fundamentos y clasificación de los recursos procesales}

\begin{note}
Imaginemos que se dictó una sentencia en un juicio y una de las partes considera que el juez se equivocó. ¿Cómo puede solicitar que un tribunal superior revise esa decisión? Ese es, precisamente, el objeto del recurso de apelación. Este capítulo recorre el camino completo: qué es un recurso, cuándo se puede interponer, cómo se concede y qué efecto produce sobre la resolución impugnada.
\end{note}

\section{Concepto de impugnación y recurso}

La voz \textit{impugnación} es un término considerablemente más amplio que el de \textit{recurso}. Según Hitter, la impugnación corresponde más a la Teoría General del Derecho, mientras que el término recurso es propio del derecho procesal. La relación entre ambos es de \textbf{género a especie}: la impugnación es el género, y el recurso es una especie de ese género.

Impugnar significa \textbf{atacar}. Pueden existir actos de impugnación fuera del proceso judicial y, dentro del proceso, existen distintos tipos de impugnaciones que no siempre constituyen recursos.

\begin{example}{Impugnación que no es recurso}
\textbf{Impugnación de dictamen pericial.} En la etapa probatoria, si se presenta un dictamen pericial médico y se corre traslado a las partes, éstas pueden atacarlo mediante un acto procesal escrito de impugnación, fundamentando por qué se lo ataca. Ese acto es claramente una impugnación, pero no es un recurso.

\textbf{Impugnación de liquidación en ejecución.} En la etapa de ejecución del proceso, si se impugna una liquidación del capital de condena más intereses por entender que los intereses están mal calculados, ese acto también es una impugnación, pero no un recurso.
\end{example}

Solo se está ante un recurso cuando lo que se impugna es una \textbf{resolución judicial}. Por ello, los recursos son mecanismos de impugnación \textbf{exclusivamente de resoluciones judiciales}: providencias simples, resoluciones interlocutorias y sentencias definitivas (reguladas a partir del art. 160 del Código Procesal).

\section{Fundamento de los recursos: el error humano}

\textbf{Calamandrei} sostenía que el fundamento de los recursos no es ni más ni menos que el \textbf{error humano}: los jueces son humanos y pueden equivocarse; si se equivocan, tiene que existir un sistema para subsanar ese error y evitar consolidar una injusticia.

\textbf{Couture} señalaba que el término recurso tiene que ver con \textit{``volver al camino de partida''}: al recurrir, el tribunal superior vuelve a recorrer un camino ya transitado en la instancia anterior donde fue dictada la resolución impugnada, para revisarlo y, en función de esa revisión, confirmar la resolución o modificarla total o parcialmente.

\section{Clasificación general de los recursos}

A continuación se presenta la clasificación general de los recursos. Cabe aclarar que, dentro de esta clasificación, únicamente se profundizará en tres de ellos, aunque resulta importante conocer el cuadro completo.

\begin{table}[H]
\centering
\caption{Clasificación general de los recursos}
\begin{tabularx}{\textwidth}{l l X}
\toprule
\textbf{Tipo} & \textbf{Categoría} & \textbf{Ejemplos} \\
\midrule
Horizontales & Remedios & Aclaratoria, revocatoria \\
Verticales ordinarios & Apelación & Recurso de apelación (proceso ordinario) \\
Verticales extraordinarios (comunes) & Regulados por ley & Inaplicabilidad de ley (fallos plenarios), recurso extraordinario federal (Ley 48) \\
Verticales extraordinarios (excepcionales) & Creación pretoriana & Sentencia arbitraria, gravedad institucional, per saltum (ya regulado) \\
Especiales & Queja & Cuando se deniega el recurso de apelación o el extraordinario \\
Parajudiciales & Fuera del proceso & Recursos administrativos, recursos arbitrales \\
\bottomrule
\end{tabularx}
\end{table}

\subsection{Recursos horizontales: aclaratoria y revocatoria}

Son horizontales porque se interponen, tramitan y resuelven \textbf{ante el mismo juez} que dictó la resolución que se impugna. Se interpone, se funda y es ese mismo juez el que la resuelve. Se dice que, más que recursos, son \textbf{remedios}, porque no hay un nuevo tribunal que vuelva a recorrer un camino de instancia anterior.

\begin{important}
Una característica fundamental diferencia a los recursos horizontales del recurso de apelación: en los horizontales se interponen y fundan \textbf{en el mismo acto}. Esto no sucede así en el recurso de apelación, donde existen dos momentos diferenciados.
\end{important}

\subsection{Recursos verticales ordinarios: la apelación}

Es vertical porque se interpone ante el juez de primera instancia (el que dictó la resolución) para que sea la \textbf{cámara} —que está por encima— la que resuelva. Tiene \textbf{dos momentos diferenciados}: la interposición y la fundamentación, que se desarrollan en profundidad en los capítulos siguientes.

\chapter{El recurso de apelación: concepto, procedencia y forma}

\section{Definición del recurso de apelación}

\begin{definition}{Recurso de apelación}
El recurso de apelación es un acto procesal de impugnación que interponen las partes o terceros ante el mismo juez que dictó una resolución —que debe ser una providencia simple, interlocutoria o definitiva— en tanto y en cuanto cause gravamen irreparable, con el objetivo de que el tribunal que está por encima de ese juzgado la revoque total o parcialmente.
\end{definition}

\section{Legitimados para interponer el recurso}

No son solo las partes en sentido estricto (actor o demandado). Cualquier sujeto procesal con \textbf{interés acreditado} puede interponer el recurso, por ejemplo:

\begin{itemize}
\item El Ministerio Público, si está representando los intereses de un menor o un incapaz.
\item Un perito que considera exiguos sus honorarios regulados.
\item El abogado interviniente respecto de la regulación de sus honorarios.
\end{itemize}

El requisito es que exista un \textbf{gravamen irreparable}: concreto, actual, que no se pueda reparar luego con el dictado de la sentencia definitiva.

\section{Resoluciones apelables: la regla de la apelabilidad}

En el proceso ordinario, la regla es la \textbf{apelabilidad}: todas las resoluciones son apelables, salvo que el Código establezca lo contrario. Esto difiere del proceso ejecutivo o del proceso sumarísimo, donde la apelabilidad está muy restringida.

\begin{important}
\textbf{Inapelable} no es lo mismo que \textbf{irrecurrible}. Si el Código dice ``irrecurrible'', esa resolución no resiste ningún tipo de recurso (ni siquiera revocatoria). Cuando dice ``inapelable'', se refiere únicamente a que no resiste el recurso de apelación; otros recursos podrían interponerse.
\end{important}

\subsection{Artículo 379: inapelabilidad de resoluciones sobre prueba}

\begin{formula}{Artículo 379 CPCCN}
Serán inapelables las resoluciones del juez sobre producción, denegación y sustanciación de las pruebas. Si se hubiere negado alguna medida, la parte interesada podrá solicitar a la cámara que la diligencie cuando el expediente le fuere remitido para que conozca del recurso contra la sentencia definitiva.
\end{formula}

\begin{example}{Denegación de prueba testimonial en audiencia preliminar}
En la audiencia preliminar, el juez decide denegar la producción de la prueba testimonial por considerarla innecesaria en función de los hechos controvertidos. La parte actora o demandada no puede apelar esa resolución porque el art. 379 la declara inapelable. Sin embargo, podrá replantear la cuestión ante la cámara cuando el expediente sea elevado para resolver el recurso contra la sentencia definitiva.
\end{example}

\subsection{Artículo 242: límite por monto}

\begin{formula}{Artículo 242 CPCCN, actualizado por Acordada 41/2019 (vigente desde el 1\textsuperscript{o} de enero de 2020)}
Serán inapelables las sentencias definitivas y las demás resoluciones, cualquiera fuere su naturaleza, que se dicten en los procesos en los que el valor cuestionado no supere los \$~300.000 (pesos trescientos mil). Esta inapelabilidad no rige en materia de alimentos ni respecto de los honorarios regulados en ese proceso.
\end{formula}

Cabe advertir que el monto se actualiza mediante acordadas de la Cámara, siendo la última la Acordada 41/2019. Se ha cuestionado la constitucionalidad de este artículo por limitar el acceso a la segunda instancia; sin embargo, la Corte Suprema ha interpretado que la doble instancia como principio constitucional-convencional (Pacto de San José de Costa Rica) rige solo en materia penal. En materia civil, el legislador puede establecer este límite sin que resulte inconstitucional ni anticonvencional. La finalidad perseguida es que los juicios de menor trascendencia económica tramiten en una sola instancia, para ser resueltos con mayor celeridad.

\section{Revocatoria con apelación en subsidio}

La única combinación posible de recursos es la \textbf{revocatoria con apelación en subsidio}, y solo procede ante una \textbf{providencia simple que cause gravamen irreparable}. En ese supuesto, el recurrente interpone la revocatoria y, para el caso de que el juez de primera instancia la rechace, subsidiariamente queda planteada la apelación para acudir a la cámara.

\begin{example}{Providencia simple que causa gravamen irreparable}
Si se contestó la demanda en término, pero el juzgado la tuvo por extemporánea por haber computado mal el plazo de 15 días (sin tener en cuenta un feriado), eso causa un gravamen irreparable. En ese supuesto: (a) puede interponerse una revocatoria con apelación en subsidio, o (b) directamente un recurso de apelación.
\end{example}

\section{Errores que comprende el recurso de apelación}

\begin{formula}{Artículo 253 CPCCN}
El recurso de apelación comprende tanto los errores in iudicando (error en el juzgamiento: fijación de los hechos o aplicación del derecho) como los errores in procedendo (error en el procedimiento de construcción de la sentencia). Por ello, en el proceso ordinario no existe un recurso de nulidad separado: el recurso de nulidad está comprendido dentro del recurso de apelación.
\end{formula}

\section{Plazo y forma de interposición}

\begin{table}[H]
\centering
\caption{Plazos de interposición según el tipo de proceso}
\begin{tabularx}{\textwidth}{X l X}
\toprule
\textbf{Proceso} & \textbf{Plazo} & \textbf{Observación} \\
\midrule
Proceso ordinario & 5 días & Desde la notificación (cédula o ministerio de ley) \\
Proceso sumarísimo & 3 días & Plazo especial previsto por el código \\
Verbal (en audiencia) & Inmediato & Se deja constancia en el acta \\
\bottomrule
\end{tabularx}
\end{table}

\section{Los dos momentos del recurso de apelación}

A diferencia de los recursos horizontales, el recurso de apelación tiene \textbf{dos momentos diferenciados}:

\begin{table}[H]
\centering
\caption{Los dos momentos del recurso de apelación}
\begin{tabularx}{\textwidth}{X X}
\toprule
\textbf{1\textsuperscript{er} momento --- Interposición} & \textbf{2\textsuperscript{do} momento --- Fundamentación} \\
\midrule
Dentro de los 5 días de notificada la resolución. Solo se declara que se interpone el recurso de apelación contra la resolución de fecha tal, por causarle gravamen irreparable. Nada más. &
Una vez concedido el recurso, en función de cómo fue concedido, se hace una crítica concreta y razonada de la resolución, explicándole a la cámara qué debería modificarse total o parcialmente y dónde está el error del juez de primera instancia. Si se funda la apelación contra la sentencia definitiva $\rightarrow$ expresión de agravios. Si se funda la apelación contra otras resoluciones $\rightarrow$ memorial. \\
\bottomrule
\end{tabularx}
\end{table}

Si no se fundamenta en tiempo y forma, se tiene por \textbf{desierto el recurso}: como si no se hubiera interpuesto, y la resolución queda firme.

\chapter{Concesión del recurso: forma, trámite y efectos}

\section{Las tres consecuencias de la concesión}

Al interponerse el recurso de apelación y ser concedido, se generan tres consecuencias:

\begin{table}[H]
\centering
\caption{Las tres consecuencias de la concesión del recurso}
\begin{tabularx}{\textwidth}{X X X}
\toprule
\textbf{Forma} & \textbf{Trámite} & \textbf{Efecto} \\
\midrule
Libre (amplia) vs. en relación (restringida) & Inmediato (regla general) vs. diferido (excepción legal) & Suspensivo (regla general) vs. no suspensivo (devolutivo --- excepción) \\
\bottomrule
\end{tabularx}
\end{table}

\section{Forma de concesión: libre vs. en relación}

\begin{table}[H]
\centering
\caption{Comparación entre concesión libre y concesión en relación}
\begin{tabularx}{\textwidth}{l X X}
\toprule
 & \textbf{Concesión libre} & \textbf{Concesión en relación} \\
\midrule
Supuesto & Solo cuando se apela la sentencia definitiva de un proceso ordinario & Todos los demás casos \\
Plazo para fundar & 10 días (expresión de agravios) & 5 días (memorial) \\
Dónde se funda & Ante la cámara & Ante el juzgado de primera instancia \\
Actividad adicional & Sí: art. 260 (5 días paralelos) & No \\
Traslado de fundamentos & 10 días (notif. por ministerio de ley) & 5 días (notif. por ministerio de ley) \\
Efecto & Siempre suspensivo & Depende de cada caso \\
Trámite & Siempre inmediato & Depende de cada caso \\
\bottomrule
\end{tabularx}
\end{table}

\subsection{Procedimiento cuando se concede en relación}

\begin{enumerate}
\item El recurrente queda notificado de la resolución (por cédula o por ministerio de ley).
\item Dentro de los 5 días interpone el recurso de apelación.
\item El juez dicta una resolución concediendo en relación el recurso (se notifica por ministerio de ley).
\item Desde que queda notificada esa concesión, el recurrente tiene 5 días para presentar el memorial (fundamentación).
\item Del memorial se corre traslado a la otra parte por 5 días (notif. por ministerio de ley).
\item Vencido ese plazo o contestado el traslado, el expediente se eleva a la cámara.
\item La cámara --- la sala correspondiente --- resuelve directamente el recurso.
\end{enumerate}

\begin{note}
\textbf{Principio de prevención.} Si a lo largo del proceso el expediente ya fue elevado a la cámara para resolver otro recurso, se remite a la misma sala que ya intervino. Si es la primera vez, se sortea la sala.
\end{note}

\subsection{Procedimiento cuando se concede libremente}

\begin{enumerate}
\item Se dicta la sentencia definitiva en un proceso ordinario.
\item Se notifica por cédula. Dentro de los 5 días se interpone el recurso.
\item El juez concede libremente todos los recursos interpuestos y el expediente es elevado directamente a la cámara.
\item La sala dicta una resolución que dice (literalmente): ``Pónganse los autos a los fines de los artículos 259 y 260 del Código Procesal. Notifíquese.''
\item Las partes son notificadas por cédula de esa resolución. Desde esa notificación corren en paralelo dos plazos: 10 días (art. 259) para expresar agravios, y 5 días (art. 260) para realizar la actividad adicional en cámara.
\item De la expresión de agravios se corre traslado a la otra parte por 10 días (notif. por ministerio de ley).
\item El tribunal dicta resolución de autos para sentencia y vota para resolver el recurso.
\end{enumerate}

\begin{formula}{Artículo 260 CPCCN --- facultades de las partes en segunda instancia (concesión libre)}
Dentro de los cinco días de la notificación de la providencia que dispone la elevación, los litigantes podrán: 1.º fundar los recursos concedidos con trámite diferido; 2.º replantear la prueba denegada o declarada negligente; 3.º acompañar documentos de fecha posterior o desconocidos al inicio del proceso; 4.º invocar hechos nuevos; 5.º pedir la absolución de posiciones sobre hechos no debatidos en la confesional de primera instancia; 6.º pedir apertura a prueba para la producción del replanteo o de los hechos nuevos.
\end{formula}

\begin{example}{Replanteo de prueba (art. 260 en relación con el art. 379)}
En la audiencia preliminar el juez denegó la prueba testimonial por considerarla innecesaria. Luego dictó sentencia rechazando la demanda por no tener por acreditada la materialidad del accidente. Si bien en su momento esa resolución era inapelable (art. 379), ahora --- al estar el expediente en cámara --- el recurrente puede replantear la prueba denegada, fundando cómo la declaración de esos testigos presenciales podría acreditar la existencia del accidente y cambiar el resultado de la sentencia.
\end{example}

\section{Límites de la alzada al resolver}

La cámara tiene tres límites para resolver:

\begin{enumerate}
\item \textbf{Expresión de agravios / memorial}: la cámara solo entiende en aquello que haya sido materia de crítica concreta y razonada. Lo que no fue objeto de agravios queda firme.
\item \textbf{Principio de congruencia}: el tribunal no puede considerar hechos que no hayan sido invocados en la demanda, contestación, reconvención o contestación de reconvención.
\item \textbf{Reformatio in peius}: la cámara no puede empeorar la situación del recurrente. Solo puede modificar en lo que fue impugnado.
\end{enumerate}

\chapter{Trámite y efectos del recurso de apelación}

\section{Trámite: inmediato vs. diferido}

La regla general es que el recurso se concede con \textbf{trámite inmediato}: una vez interpuesto y concedido, de modo continuo se computa el plazo para fundar, se corre traslado y la cámara resuelve. La excepción es el \textbf{trámite diferido}, cuando la ley lo dispone: el momento de fundamentación, sustanciación y resolución se difiere para cuando el expediente es elevado a la cámara por el recurso contra la sentencia definitiva.

\begin{important}
Advertencia terminológica: el legislador utiliza la expresión ``efecto diferido'' en el código, pero en realidad se trata de una cuestión de trámite, no de efecto. Cuando el código dice ``efecto diferido'' debe leerse ``trámite diferido''.
\end{important}

\subsection{Supuestos de trámite diferido}

\begin{formula}{Artículo 366 --- hechos nuevos}
La resolución que rechaza el planteo de hechos nuevos puede apelarse con trámite diferido. La fundamentación se hace cuando el expediente llega a la cámara para resolver el recurso contra la sentencia definitiva. Finalidad: evitar que el expediente sea elevado varias veces antes de la sentencia.
\end{formula}

\begin{formula}{Artículo 69 --- costas y honorarios incidentales}
Las apelaciones sobre costas y honorarios regulados a lo largo del proceso se conceden con trámite diferido para el mismo momento. Misma finalidad: celeridad y economía procesal.
\end{formula}

\section{Efectos: suspensivo vs. no suspensivo (devolutivo)}

El efecto hace referencia a qué sucede con la resolución impugnada: ¿se suspenden sus efectos o no? La regla general es el \textbf{efecto suspensivo}: al interponerse el recurso, se suspende la ejecutoriedad de la resolución hasta que la cámara resuelva. La razón es que la cámara puede modificar o revocar la resolución, por lo que tiene sentido no ejecutarla mientras está siendo revisada.

\subsection{Origen histórico del término ``devolutivo''}

El término ``devolutivo'' que emplea el código tiene origen en el derecho romano: el único con jurisdicción era el emperador, quien la delegaba en los jueces. Cuando una parte impugnaba una decisión, se ``devolvía'' la jurisdicción al emperador. En el sistema actual, todos los jueces (desde el de primera instancia hasta la Corte Suprema) tienen la misma función jurisdiccional; lo que varía es el ámbito de su competencia.

\begin{important}
El término ``devolutivo'' puede llevar a confusión: cuando el código dice ``devolutivo'', debe leerse ``no suspensivo''.
\end{important}

\subsection{Supuestos de efecto no suspensivo}

\begin{example}{Alimentos, medidas cautelares y beneficio de litigar sin gastos}
\textbf{Alimentos (sentencia que admite).} La sentencia que fija alimentos a cargo del demandado puede ejecutarse aunque esté pendiente el recurso de apelación del demandado ante la cámara.

\textbf{Medidas cautelares (art. 198).} La resolución que concede una medida cautelar (por ejemplo, un embargo) puede apelarse por el afectado, pero el recurso tiene efecto no suspensivo: el embargo se traba y continúa trabándose sin importar que el recurso esté pendiente ante la cámara.

\textbf{Beneficio de litigar sin gastos (art. 81).} La resolución que concede el beneficio tiene efecto no suspensivo: se aplica de inmediato, aunque haya recurso en trámite.
\end{example}

\section{Artículo 250: procedimiento según el tipo de resolución}

\begin{table}[H]
\centering
\caption{Procedimiento según el tipo de resolución (efecto no suspensivo)}
\begin{tabularx}{\textwidth}{X X}
\toprule
\textbf{Resolución es sentencia} & \textbf{Resolución es interlocutoria} \\
\midrule
El expediente se eleva a la cámara para que pueda revisar todo el proceso y resolver el recurso. Para ejecutar lo resuelto en la sentencia, se forma un incidente de ejecución de sentencia que queda en primera instancia (con copia certificada de la sentencia). Expediente principal $\rightarrow$ cámara. Incidente de ejecución $\rightarrow$ primera instancia. &
El expediente queda en primera instancia para continuar el proceso (la sentencia definitiva aún no fue dictada). Se forma un incidente de apelación con las copias indicadas por el juez, que se eleva a la cámara. El apelante tiene 5 días para acompañar esas copias. Expediente principal $\rightarrow$ primera instancia. Incidente de apelación $\rightarrow$ cámara. \\
\bottomrule
\end{tabularx}
\end{table}

Si el apelante no acompaña las copias para formar el incidente dentro de los 5 días, se declara la \textbf{deserción} del recurso (como si no se hubiera interpuesto).

\chapter{Las seis reglas del recurso de apelación}

A continuación se sintetizan todas las reglas estudiadas respecto del recurso de apelación.

\begin{longtable}{
    >{\bfseries\raggedright\arraybackslash}p{1.4cm}
    >{\bfseries\raggedright\arraybackslash}p{3.4cm}
    >{\raggedright\arraybackslash}p{9.2cm}
}
\caption{Las seis reglas del recurso de apelación} \\
\toprule
\textbf{N.\textsuperscript{o}} & \textbf{Regla} & \textbf{Contenido} \\
\midrule
\endfirsthead
\toprule
\textbf{N.\textsuperscript{o}} & \textbf{Regla} & \textbf{Contenido} \\
\midrule
\endhead
\bottomrule
\endfoot
Regla 1 & Alcance del recurso & El recurso de apelación comprende tanto los errores in iudicando (fijación de hechos / aplicación del derecho) como los errores in procedendo (procedimiento de construcción de la sentencia). El recurso de nulidad está incluido dentro del de apelación (art. 253 CPCCN). \\
Regla 2 & Plazo de interposición & El plazo para interponer el recurso es de 5 días, salvo que la ley disponga lo contrario (por ejemplo, proceso sumarísimo: 3 días). \\
Regla 3 & Forma de concesión & El recurso se concede siempre en relación, salvo un único supuesto en que se concede libremente: cuando lo que se apela es la sentencia definitiva dictada en un proceso ordinario. \\
Regla 4 & Libre $\rightarrow$ siempre suspensivo e inmediato & Cuando el recurso se concede libremente (sentencia definitiva de proceso ordinario), siempre es con efecto suspensivo y trámite inmediato. \\
Regla 5 & Trámite & El recurso se concede siempre con trámite inmediato, salvo que la ley disponga expresamente el trámite diferido (por ejemplo, arts. 366 y 69 CPCCN). \\
Regla 6 & Efecto & El recurso se concede siempre con efecto suspensivo, salvo que la ley disponga lo contrario (efecto no suspensivo / devolutivo, por ejemplo alimentos, medidas cautelares, beneficio de litigar sin gastos). \\
\end{longtable}

\chapter{Ejercicios prácticos y mecanismos especiales}

\section{Admisibilidad vs. fundabilidad: ¿quién analiza qué?}

\begin{exercise}{Admisibilidad y fundabilidad}
¿Cuál va a ser el juez o tribunal que analiza, en el caso del recurso de apelación, los requisitos de admisibilidad del recurso y los de fundabilidad?
\end{exercise}

\textbf{Resolución.} \textbf{Admisibilidad $\rightarrow$ juez de primera instancia}: analiza si el recurso fue interpuesto en el plazo legal, si la resolución es susceptible de apelación y si quien lo interpone es el legitimado. Eso lo hace al momento de conceder o denegar el recurso. \textbf{Fundabilidad $\rightarrow$ cámara}: analiza si la expresión de agravios o el memorial constituyen una crítica concreta y razonada, y decide si hace lugar o no al recurso (confirma, modifica o revoca la resolución).

\section{Caso práctico: contestación de demanda declarada extemporánea}

\begin{exercise}{Contestación de demanda declarada extemporánea}
Se trata de los abogados del demandado en un proceso ordinario de daños y perjuicios por mala praxis médica. Tienen 15 días para contestar demanda y la contestan el día 16, dentro del plazo de gracia (las dos primeras horas del día, antes de las 9:30~a.\,m.). El día siguiente, el juzgado emite esta resolución: ``Buenos Aires, 22 de junio de 2020. DESGLÓSESE por extemporáneo.'' Firma: María Eugenia Gómez, Secretaria. ¿Qué corresponde hacer? (El juzgado no tuvo en cuenta un feriado intermedio y por eso computó mal el plazo.)
\end{exercise}

Se analizan a continuación posibles respuestas:

\begin{itemize}
\item Interponer una revocatoria (art. 238 y ss.) $\rightarrow$ respuesta incorrecta.
\item Interponer revocatoria con apelación en subsidio $\rightarrow$ respuesta incorrecta.
\item Interponer directamente un recurso de apelación $\rightarrow$ respuesta incorrecta.
\end{itemize}

\begin{important}
Todas las respuestas anteriores son incorrectas. ¿Por qué? Porque la resolución no está firmada por el juez, sino por la secretaria del juzgado. El secretario y el prosecretario actúan por delegación. Solo el juez tiene jurisdicción; solo el juez puede revocar por contrario imperio. La resolución no es susceptible ni de revocatoria ni de apelación. La vía correcta es la \textbf{revocación nominada del artículo 38 ter del CPCCN}.
\end{important}

\section{Revocación nominada: artículo 38 ter CPCCN}

\begin{formula}{Artículo 38 ter CPCCN}
Dentro del plazo de tres días desde que se notifican de la resolución dictada por el secretario o prosecretario administrativo, las partes pueden requerir al juez que deje sin efecto lo dispuesto por el secretario o prosecretario administrativo.
\end{formula}

A continuación se señalan las diferencias con la revocatoria del art. 238:

\begin{table}[H]
\centering
\caption{Revocatoria (art. 238) vs. revocación nominada (art. 38 ter)}
\begin{tabularx}{\textwidth}{l X X}
\toprule
 & \textbf{Revocatoria (art. 238)} & \textbf{Revocación nominada (art. 38 ter)} \\
\midrule
Plazo & 5 días & 3 días \\
Momento de fundamentación & Simultáneo a la interposición & Simultáneo a la interposición \\
Resolución que ataca & Providencias simples dictadas por el juez & Resoluciones dictadas por secretario o prosecretario \\
Quién resuelve & El mismo juez & El juez (revisa lo dispuesto por el auxiliar) \\
\bottomrule
\end{tabularx}
\end{table}

\section{Reposición in extremis: creación pretoriana}

Se trata de un mecanismo impugnativo \textbf{no regulado en el código}, creado por la jurisprudencia. No debe confundirse ni con la revocatoria del art. 238 ni con la revocación del art. 38 ter.

\begin{important}
\textbf{Requisitos de la reposición in extremis:}
\begin{itemize}
\item Se trata de una resolución que NO es una providencia simple.
\item No existe otro remedio ni recurso posible.
\item El error es de índole sustancial (de fondo), no meramente formal.
\item El error es notorio: de mantenerse, se consolidaría una injusticia muy grande para alguno de los litigantes.
\item Se interpone y se funda en el mismo acto. Plazo: 3 días.
\end{itemize}
\end{important}

\begin{example}{Sentencia inapelable por monto con error sustancial}
Una sentencia definitiva dictada en un proceso donde la demanda es de \$~200.000 (menor al mínimo del art. 242) es inapelable. Si esa sentencia tiene un error sustancial en la fijación de los hechos que constituye una injusticia notoria, la parte afectada puede interponer la reposición in extremis dentro de los 3 días, fundándola en ese mismo acto.
\end{example}

\subsection{Jurisprudencia sobre reposición in extremis}

\begin{formula}{Cámara Civil, Sala E --- sumario citado}
Las resoluciones de este tribunal de alzada no son en principio susceptibles de recurso de reposición, por no revestir aquéllas el carácter de providencias simples. Sin embargo, este principio reconoce excepciones en circunstancias excepcionales cuando, por mediar un error esencial en la apreciación de los antecedentes del caso, el mantenimiento de la situación conduciría a un resultado injusto que es deber de los jueces modificar, para preservar a los litigantes.
\end{formula}

Cabe agregar que la Corte Suprema de Justicia también ha sostenido que sus decisiones no son susceptibles de recurso de reposición, aunque ese principio puede reconocer excepciones cuando se trata de situaciones serias e inequívocas que demuestran con nitidez manifiesta el error que se pretende subsanar.

% TODO: contenido incompleto o ambiguo en las notas originales — no se identifica en el material la cita completa (tribunal, sala, fecha o número de fallo) de la jurisprudencia de la Corte Suprema mencionada en este párrafo.

\chapter{Cuadro Cornell de repaso}

\section*{Recurso de Apelación --- Derecho Procesal Civil y Comercial}
\addcontentsline{toc}{section}{Recurso de Apelación --- Derecho Procesal Civil y Comercial}

\begin{longtable}{
    >{\raggedright\arraybackslash}p{0.30\textwidth}
    >{\raggedright\arraybackslash}p{0.62\textwidth}
}
\toprule
\textbf{Preguntas / palabras clave} & \textbf{Notas / respuestas} \\
\midrule
\endfirsthead
\toprule
\textbf{Preguntas / palabras clave} & \textbf{Notas / respuestas} \\
\midrule
\endhead
\bottomrule
\endfoot

\textbf{¿Qué diferencia hay entre impugnación y recurso?} &
La impugnación es el género (todo acto de ataque contra un acto jurídico, dentro o fuera del proceso). El recurso es la especie: solo ataca resoluciones judiciales. No toda impugnación es un recurso, pero todo recurso es una impugnación. \\

\textbf{¿Cuál es el fundamento de los recursos?} &
El error humano (Calamandrei). Los jueces pueden equivocarse; los recursos son el sistema para subsanar esos errores y evitar que se consolide una injusticia. Couture agrega: ``recurrir'' significa volver al camino para que un tribunal superior lo revise. \\

\textbf{¿Cuándo se concede libremente?} &
Solo cuando se apela la sentencia definitiva de un proceso ordinario. En todos los demás casos, el recurso se concede en relación. \\

\textbf{¿Qué son los dos momentos del recurso de apelación?} &
1.\textsuperscript{o} Interposición: dentro de los 5 días de notificada la resolución. Solo se declara que se interpone el recurso. 2.\textsuperscript{o} Fundamentación: una vez concedido el recurso, se hace la crítica concreta y razonada (expresión de agravios o memorial). Si no se fundamenta, el recurso se declara desierto. \\

\textbf{¿Qué es la expresión de agravios?} &
La fundamentación del recurso cuando se apela la sentencia definitiva. Es una crítica concreta y razonada que le explica a la cámara dónde se equivocó el juez (en los hechos, en el derecho o en el procedimiento). Fija el límite del conocimiento de la alzada. \\

\textbf{¿Qué es el efecto suspensivo?} &
Al interponerse el recurso, se suspende la ejecutoriedad de la resolución impugnada hasta que la cámara resuelva. Es la regla general. Excepción: efecto no suspensivo (devolutivo) cuando la ley lo dispone expresamente (alimentos, medidas cautelares, beneficio de litigar sin gastos). \\

\textbf{¿Qué es el trámite diferido?} &
El momento de fundamentación, sustanciación y resolución del recurso se difiere para cuando el expediente sea elevado a la cámara por el recurso contra la sentencia definitiva. Supuestos: art. 366 (hechos nuevos) y art. 69 (costas y honorarios incidentales). Finalidad: celeridad y economía procesal. \\

\textbf{¿Qué es la reformatio in peius?} &
El tribunal de alzada no puede empeorar la situación del recurrente. Solo puede modificar la resolución en lo que fue materia de agravios. Es uno de los tres límites de la alzada (junto con la expresión de agravios y el principio de congruencia). \\

\textbf{¿Cuándo procede la reposición in extremis?} &
Cuando: (1) no es una providencia simple, (2) no existe otro remedio o recurso posible, (3) el error es sustancial (no formal), (4) el error es notorio y su mantenimiento consagraría una injusticia grave. Se interpone y funda en el mismo acto (3 días). Es una creación pretoriana; no está en el código. \\

\textbf{¿Cuál es la diferencia entre el art. 38 ter y la revocatoria (art. 238)?} &
Art. 238 (revocatoria): plazo de 5 días, ataca providencias simples dictadas por el juez. Art. 38 ter (revocación nominada): plazo de 3 días, ataca resoluciones del secretario o prosecretario. En ambos, se interpone y funda en el mismo acto. \\

\textbf{¿Qué sucede si no se acompañan las copias en 5 días (efecto no suspensivo en interlocutoria)?} &
Se declara la deserción del recurso: es como si no se hubiera interpuesto, y la resolución queda firme para esa parte. \\

\textbf{¿Qué analiza el juez de primera instancia vs. la cámara?} &
Primera instancia: admisibilidad (plazo, resolución susceptible, legitimado, gravamen irreparable, monto). Cámara: fundabilidad (si la expresión de agravios o memorial constituye crítica concreta y razonada) y el fondo del recurso. \\

\midrule

\multicolumn{2}{p{0.92\textwidth}}{
\textbf{Resumen:} El recurso de apelación es el mecanismo central de impugnación de resoluciones judiciales en el proceso civil. Su fundamento radica en la falibilidad humana: los jueces pueden equivocarse, y el sistema procesal debe ofrecer una vía para que un tribunal superior revise esas decisiones. A diferencia de los recursos horizontales (aclaratoria y revocatoria), que se interponen, tramitan y resuelven ante el mismo juez, el recurso de apelación es vertical: se interpone ante el juez que dictó la resolución, pero lo resuelve la cámara que está por encima. El recurso comprende tanto errores de juzgamiento (in iudicando: fijación de hechos o aplicación del derecho) como errores de procedimiento (in procedendo), lo que lo convierte en un mecanismo integral que absorbe al recurso de nulidad. Tiene dos momentos diferenciados: interposición (5 días) y fundamentación (expresión de agravios o memorial), que no solo hace viable el recurso, sino que fija los límites del conocimiento de la alzada. La cámara solo entenderá en lo que fue objeto de crítica concreta y razonada, respetando además el principio de congruencia y la prohibición de la reformatio in peius. La concesión del recurso genera tres consecuencias: una forma (libre --- solo para sentencias definitivas de procesos ordinarios --- vs. en relación para todos los demás casos), un trámite (inmediato como regla, diferido como excepción legal para hechos nuevos y costas incidentales) y un efecto (suspensivo como regla, no suspensivo/devolutivo como excepción para alimentos, medidas cautelares y beneficio de litigar sin gastos). Estas tres dimensiones interactúan: la concesión libre siempre es con efecto suspensivo y trámite inmediato. Dos mecanismos especiales completan el panorama: la revocación nominada (art. 38 ter), que en 3 días y con fundamentación simultánea ataca resoluciones del secretario o prosecretario --- no susceptibles de revocatoria ni apelación ---; y la reposición in extremis, creación pretoriana que opera cuando no existe otro remedio, el error es sustancial y notorio, y su mantenimiento implicaría una injusticia grave. Identificar correctamente el tipo de resolución y quién la firmó es la clave para elegir la vía impugnativa correcta.
} \\

\end{longtable}

\end{document}
